\section*{Chapitre \ref{ch:b1:respiration-fermentation} --- Respiration cellulaire et fermentation}

\begin{solution}{exo:b1:respiration-fermentation:1}
\omterm{def:b1:respiration-fermentation:glycolysis}{Glycolyse}, \omterm{def:b1:eukaryotic-cell:organelle}{cytosol}~: 2 pyruvates, 2 \omterm{def:b1:water-small-molecules:atp}{ATP}, 2 NADH. Oxydation du pyruvate
et \omterm{def:b1:respiration-fermentation:krebs}{cycle de Krebs}, \omterm{def:b1:eukaryotic-cell:mitochondrion}{matrice mitochondriale}~: 6 $\mathrm{CO_2}$, 8 NADH, 2
$\mathrm{FADH_2}$, 2 GTP. Phosphorylation oxydative, membrane interne~:
de l'eau à partir de l'oxygène, environ 26 \omterm{def:b1:water-small-molecules:atp}{ATP}.
\end{solution}

\begin{solution}{exo:b1:respiration-fermentation:2}
Glucose $+ 2\,\mathrm{NAD^+} + 2\,\mathrm{ADP} + 2\,\mathrm{P_i} \to
2$ pyruvate $+ 2\,\mathrm{NADH} + 2\,\mathrm{H^+} + 2\,\mathrm{ATP} +
2\,\mathrm{H_2O}$. \omterm{def:b1:respiration-fermentation:glycolysis}{Phosphorylation au niveau du substrat}~: un phosphate
est transféré directement d'un intermédiaire riche en énergie
(1,3-bisphosphoglycérate, phosphoénolpyruvate) à l'ADP par une \omterm{def:b1:enzymes:enzyme}{enzyme},
sans gradient membranaire ni oxygène.
\end{solution}

\begin{solution}{exo:b1:respiration-fermentation:3}
NADH et $\mathrm{CO_2}$ (isocitrate en $\alpha$-cétoglutarate), NADH
et $\mathrm{CO_2}$ ($\alpha$-cétoglutarate en succinyl-CoA), GTP
(succinyl-CoA en succinate), $\mathrm{FADH_2}$ (succinate en fumarate),
NADH (malate en oxaloacétate). Les deux $\mathrm{CO_2}$ viennent des
deux décarboxylations des acides à six et à cinq carbones.
\end{solution}

\begin{solution}{exo:b1:respiration-fermentation:4}
Elle régénère le $\mathrm{NAD^+}$ pour que la \omterm{def:b1:respiration-fermentation:glycolysis}{glycolyse} puisse continuer
et rapporter ses 2 \omterm{def:b1:water-small-molecules:atp}{ATP} par glucose. Elle n'oxyde pas le carbone,
n'extrait pas l'énergie qui reste (le produit la conserve presque tout
entière) et ne fabrique pas plus d'un quinzième de l'\omterm{def:b1:water-small-molecules:atp}{ATP} de la
respiration.
\end{solution}

\begin{solution}{exo:b1:respiration-fermentation:5}
Vers l'oxygène~: $-2\times 96\,500\times 1.14 = \qty{-220}{kJ/mol}$,
soit quatre \omterm{def:b1:water-small-molecules:atp}{ATP} à \qty{50}{kJ}. Vers l'ubiquinone~: $-2\times 96\,500\times 0.37 =
\qty{-71}{kJ/mol}$, un \omterm{def:b1:water-small-molecules:atp}{ATP}.
\end{solution}

\begin{solution}{exo:b1:respiration-fermentation:6}
Muscle~: $2 + 2$ (GTP) $+ 8\times 2.5$ (NADH mitochondrial) $+ 2\times
1.5$ ($\mathrm{FADH_2}$) $+ 2\times 1.5$ (NADH cytosolique par le \omterm{def:b1:respiration-fermentation:carriers}{FAD}) $=
30$. Foie~: le NADH cytosolique donne $2\times 2.5$~: 32.
\end{solution}

\begin{solution}{exo:b1:respiration-fermentation:7}
QR $= 2.2$. Glucose respiré $x$~: $6x = 1.0$, $x = \qty{0.167}{mmol}$~;
$\mathrm{CO_2}$~: $6x + 2y = 2.2$, $2y = 1.2$, $y = \qty{0.6}{mmol}$.
Fermenté $0.6/0.767 = \qty{78}{\%}$ du glucose, respiré \qty{22}{\%}.
\end{solution}

\begin{solution}{exo:b1:respiration-fermentation:8}
La consommation d'oxygène augmente (la chaîne, libérée de la
contre-pression du gradient, tourne à plein régime), la synthèse d'\omterm{def:b1:water-small-molecules:atp}{ATP}
s'arrête (plus de gradient pour entraîner la synthase), et toute
l'énergie des électrons paraît en chaleur. Les patients brûlaient vite
leur \omterm{def:b1:lipids:triglyceride}{graisse} et mouraient d'une température corporelle que rien ne
pouvait faire baisser.
\end{solution}

\begin{solution}{exo:b1:respiration-fermentation:9}
La chaîne s'engorge~: chaque transporteur devient réduit, aucun proton
n'est pompé, aucun \omterm{def:b1:water-small-molecules:atp}{ATP} n'est fabriqué. Le NADH n'est plus réoxydé, de
sorte que le \omterm{def:b1:respiration-fermentation:krebs}{cycle de Krebs} et la pyruvate déshydrogénase s'arrêtent
faute de $\mathrm{NAD^+}$. La \omterm{def:b1:respiration-fermentation:glycolysis}{glycolyse}, libérée de l'inhibition par
l'\omterm{def:b1:water-small-molecules:atp}{ATP}, tourne vite et doit régénérer son $\mathrm{NAD^+}$ en réduisant
le pyruvate en lactate, qui inonde le sang.
\end{solution}

\begin{solution}{exo:b1:respiration-fermentation:10}
$8\times 10 + 7\times 2.5 + 7\times 1.5 - 2 = 80 + 17.5 + 10.5 - 2 =
106$ \omterm{def:b1:water-small-molecules:atp}{ATP}~; $106/16 = 6.6$ par carbone. Oxygène~: le palmitate demande 23
$\mathrm{O_2}$ pour 106 \omterm{def:b1:water-small-molecules:atp}{ATP}, soit \qty{0.22}{O_2} par \omterm{def:b1:water-small-molecules:atp}{ATP}~; le glucose 6
pour 30, soit \qty{0.20}{O_2} par \omterm{def:b1:water-small-molecules:atp}{ATP}. Par unité d'oxygène, le glucose
est \qty{10}{\%} meilleur~: le phoque, limité par l'oxygène sous l'eau,
préfère les \omterm{def:b1:carbohydrates:monosaccharide}{glucides}~; par unité de masse, la \omterm{def:b1:lipids:triglyceride}{graisse} est deux fois
meilleure~: le colibri, limité par la masse en vol, migre sur sa
\omterm{def:b1:lipids:triglyceride}{graisse}.
\end{solution}

\begin{solution}{exo:b1:respiration-fermentation:11}
\omterm{def:b1:water-small-molecules:atp}{ATP} seul~: $5/3 < \qty{2}{s}$. Phosphocréatine~: $25/3 = \qty{8}{s}$.
\omterm{def:b1:respiration-fermentation:fermentation}{Fermentation}~: $2\times 3 = \qty{6}{mmol/L}$ d'\omterm{def:b1:water-small-molecules:atp}{ATP} par seconde, assez
pour la demande, tant que durent le \omterm{def:b1:carbohydrates:polysaccharide}{glycogène} et la tolérance au
lactate (quelques dizaines de secondes). La livraison d'oxygène met une
minute ou plus à monter et ne fournit au plus qu'environ
\qty{1}{mmol/L} d'\omterm{def:b1:water-small-molecules:atp}{ATP} par seconde dans une fibre~: une course de dix
secondes est finie avant que la respiration n'ait commencé à aider, et
elle est payée par la phosphocréatine et la \omterm{def:b1:respiration-fermentation:fermentation}{fermentation}.
\end{solution}

\begin{solution}{exo:b1:respiration-fermentation:12}
Le gradient stocke de l'énergie sous forme d'une différence de
concentration en protons et de potentiel électrique de part et d'autre
d'une membrane isolante --- exactement un condensateur chargé doublé
d'une pile de concentration~; la chaîne est le chargeur, entraîné par la
chute des électrons~; la synthase est un moteur tourné par le courant de
protons, exactement une turbine, qui couple le flux à une rotation
mécanique qui fabrique de l'\omterm{def:b1:water-small-molecules:atp}{ATP}. L'analogie est approximative en ceci
que la «~pile~» ne contient que quelques millisecondes de la
consommation de la \omterm{def:b1:cell-unit-of-life:cell}{cellule} et doit être rechargée en permanence, que les
autres transporteurs de la même membrane y puisent aussi, et que la
turbine peut tourner à l'envers, hydrolysant de l'\omterm{def:b1:water-small-molecules:atp}{ATP} pour pomper des
protons quand la chaîne défaille.
\end{solution}

\begin{solution}{pb:b1:respiration-fermentation:1}
\textbf{1.} \omterm{def:b1:respiration-fermentation:glycolysis}{Glycolyse}~: 2 \omterm{def:b1:water-small-molecules:atp}{ATP}, 2 NADH (cytosoliques). Pyruvate
déshydrogénase~: 2 NADH. Krebs~: 6 NADH, 2 $\mathrm{FADH_2}$, 2 GTP.
\textbf{2.} NADH mitochondrial $8\times 10 = 80$~; $\mathrm{FADH_2}$
$2\times 6 = 12$~; NADH cytosolique en $\mathrm{FADH_2}$ $2\times 6 = 12$~:
104 protons.
\textbf{3.} $104/4 = 26$ \omterm{def:b1:water-small-molecules:atp}{ATP}~; total $26 + 2 + 2 = 30$.
\textbf{4.} 2 \omterm{def:b1:water-small-molecules:atp}{ATP}.
\textbf{5.} 15.
\textbf{6.} Avec air $30\times 50 = \qty{1500}{kJ}$~: \qty{52}{\%} des
2870. Sans~: \qty{100}{kJ}, soit \qty{3.5}{\%} des 2870 mais
\qty{43}{\%} des 235 libérés --- la \omterm{def:b1:respiration-fermentation:fermentation}{fermentation} capte efficacement ce
qu'elle libère~; elle en libère seulement peu.
\textbf{7.} Fiole aérée $30/30 = \qty{1}{mmol}$ de glucose par heure~;
fiole close $30/2 = \qty{15}{mmol}$.
\textbf{8.} Fiole close~: \qty{30}{mmol} d'éthanol et \qty{30}{mmol} de
$\mathrm{CO_2}$ par heure. Fiole aérée~: \qty{6}{mmol} de
$\mathrm{CO_2}$ libérées et \qty{6}{mmol} d'$\mathrm{O_2}$ consommées.
\textbf{9.} $2\times 10.5 = \qty{21}{g}$ par mole, soit \qty{0.12}{g}
par gramme de glucose.
\textbf{10.} $0.5\times 180 = \qty{90}{g}$ par mole, soit $90/10.5 =
8.6$ \omterm{def:b1:water-small-molecules:atp}{ATP} d'équivalent contre 30 disponibles~: l'\omterm{def:b1:water-small-molecules:atp}{ATP} est en excès~; le
rendement est limité par le carbone (la moitié du glucose est oxydée en
$\mathrm{CO_2}$ pour faire l'\omterm{def:b1:water-small-molecules:atp}{ATP}, et la biomasse réclame des squelettes
carbonés et du pouvoir réducteur), et l'\omterm{def:b1:water-small-molecules:atp}{ATP} excédentaire est dépensé en
entretien.
\textbf{11.} Pour obtenir le même \omterm{def:b1:water-small-molecules:atp}{ATP}, la culture close consomme quinze
fois plus de glucose (question 7) et croît cinq fois moins par gramme
(question 9)~: consommation rapide du sucre, croissance faible, alcool
--- l'observation de Pasteur. La phosphofructokinase, inhibée par l'\omterm{def:b1:water-small-molecules:atp}{ATP}
et le citrate en aération, est libérée quand la respiration s'arrête.
\textbf{12.} $\mathrm{CO_2}$ $44/44 = \qty{1.0}{mmol}$~; $\mathrm{O_2}$
$20/32 = \qty{0.625}{mmol}$~: QR $= 1.6$.
\textbf{13.} $\mathrm{O_2}$~: $6x = 0.625$, $x = \qty{0.104}{mmol}$.
$\mathrm{CO_2}$~: $6x + 2y = 1.0$, $2y = 0.375$, $y = \qty{0.1875}{mmol}$.
\textbf{14.} Fermenté $0.1875/0.292 = \qty{64}{\%}$ du glucose. \omterm{def:b1:water-small-molecules:atp}{ATP}
issu de la \omterm{def:b1:respiration-fermentation:fermentation}{fermentation} $2y = 0.375$ contre $30x = 3.12$ issus de la
respiration~: \qty{11}{\%}.
\textbf{15.} $16/23 = 0.70$.
\textbf{16.} À jeun, l'animal brûle de la \omterm{def:b1:lipids:triglyceride}{graisse} (QR 0.7)~; nourri, il
brûle les \omterm{def:b1:carbohydrates:monosaccharide}{glucides} du repas (1.0).
\textbf{17.} $96\,500\times 0.22 = \qty{21.2}{kJ/mol}$.
\textbf{18.} $4\times 21.2 = \qty{85}{kJ}$ pour un \omterm{def:b1:water-small-molecules:atp}{ATP} de \qty{50}{kJ}~:
\qty{59}{\%}.
\textbf{19.} $2\times 96\,500\times 1.14 = \qty{220}{kJ}$~; dix protons
stockent $10\times 21.2 = \qty{212}{kJ}$~: \qty{96}{\%} --- la chaîne
gaspille peu~; les pertes sont à la synthase et aux transporteurs.
\textbf{20.} $0.96\times 0.59 = 0.57$~; $2.5\times 50/220 = 0.57$~: le
même chiffre par les deux voies.
\textbf{21.} La consommation d'oxygène augmente, le rendement en \omterm{def:b1:water-small-molecules:atp}{ATP}
tombe vers celui de la \omterm{def:b1:respiration-fermentation:fermentation}{fermentation} (2 par glucose issus de la
\omterm{def:b1:respiration-fermentation:glycolysis}{glycolyse}, et la PFK étant libérée, la consommation de glucose
augmente), le dégagement de chaleur augmente, la croissance s'effondre.
\textbf{22.} La synthèse d'\omterm{def:b1:water-small-molecules:atp}{ATP} s'arrête~; le gradient ne peut plus se
décharger, la chaîne cale contre lui, la consommation d'oxygène cesse,
la \omterm{def:b1:cell-unit-of-life:cell}{cellule} fermente ce qu'elle peut. Ajouter ensuite un découplant
laisse les protons refluer~: la chaîne et la consommation d'oxygène
reprennent, toujours sans \omterm{def:b1:water-small-molecules:atp}{ATP}.
\textbf{23.} Le NADH ne fait pomper que les complexes III et IV~:
$6/4 = 1.5$~; par glucose $8\times 1.5 + 2\times 1.5 + 2\times 1.5 + 4 = 22$ \omterm{def:b1:water-small-molecules:atp}{ATP}.
\textbf{24.} Sans oxygène, la chaîne ne peut pas prendre d'électrons, de
sorte que le NADH n'est pas réoxydé et que le pool reste réduit~; les
trois déshydrogénases à \omterm{def:b1:respiration-fermentation:carriers}{NAD} du cycle n'ont plus de $\mathrm{NAD^+}$ et
le cycle s'arrête --- c'est pourquoi la \omterm{def:b1:respiration-fermentation:fermentation}{fermentation} doit régénérer le
$\mathrm{NAD^+}$ dans le \omterm{def:b1:eukaryotic-cell:organelle}{cytosol} pour que la \omterm{def:b1:respiration-fermentation:glycolysis}{glycolyse} continue.
\textbf{25.} 30 \omterm{def:b1:water-small-molecules:atp}{ATP} avec air, 2 sans~: un rapport de 15~; P/O de 2.5
pour le NADH (10 protons pour 4 par \omterm{def:b1:water-small-molecules:atp}{ATP}) et de 1.5 pour le
$\mathrm{FADH_2}$ (6 pour 4).
\end{solution}
